(1) 初始基态满足 \(\langle x^2\rangle_i=\hbar/(2m\omega_i)\)、\(\langle p^2\rangle_i=\hbar m\omega_i/2\)。由于瞬时变频不改变 \(x,p\) 的量子态，最终平均能量 \(\langle H_f\rangle=\hbar[\omega_i+\omega_f^2/\omega_i]/4\)。绝热频率变化将基态带到最终基态，能量仅 \(\hbar\omega_f/2\)。

(2) 定义 \(r=\tfrac12\ln(\omega_f/\omega_i)\)，Bogoliubov关系可写 \(a_f=\cosh r\,a_i+\sinh r\,a_i^\dagger\)（符号与压缩相位约定有关）。所以初始真空相对于新算符是压缩真空，\(\boxed{\langle n_f\rangle=\sinh^2r=(\omega_f/\omega_i+\omega_i/\omega_f-2)/4}\)，严格非负。

(3) 波函数关于 \(x\) 为偶函数，而最终本征函数宇称为 \((-1)^n\)，故所有奇数 \(n\) 的重叠均为零。偶数态概率为 \(\boxed{P_{2j}=\frac{(2j)!}{2^{2j}(j!)^2}\frac{\tanh^{2j}r}{\cosh r}}\)，可用二项级数核验归一；基态重叠概率 \(P_0=1/\cosh r=2\sqrt{\omega_i\omega_f}/(\omega_i+\omega_f)\)。

(4) 额外非绝热功为 \(W_{\rm exc}=\langle H_f\rangle-\hbar\omega_f/2=\hbar\omega_f\sinh^2r\ge0\)。压缩真空满足 \(\operatorname{Var}(n_f)=2\sinh^2r(1+\sinh^2r)\)，比相同平均数的泊松分布更宽。慢速绝热变频保持瞬时基态占据 \(n=0\)，没有这项激发能。
